%% parts/p08/many-valued-and-fuzzy.tex --- Part 8.5
%% Inherited source: v1.1 lines 13584--13787; exercise solutions 16539--16556.
%% R-037 control: saturation is defined before the finite-matrix theorem, and
%% its complete source proof is migrated to annex:p8.matrix.1.
%% OPEN source-exercise obligation: ex:many-valued-paraconsistent-01 and
%% ex:many-valued-paraconsistent-02, including their supplied feedback, remain
%% owed in full at exact source-mapped destinations.  The theorem, observation,
%% and example below are separate later material and do not discharge them.

\chapter{Many-Valued and Fuzzy Logics}
\label{sec:p08-many-valued-and-fuzzy}
\chaptermark{Many-Valued and Fuzzy Logics}

\section{Finite logical matrices}

\begin{hbdefinition}{def:p08-logical-matrix}{Finite propositional logical matrix}
Fix a propositional signature.  A \emph{finite logical matrix} is a structure
\[
  \mathcal M=(\mathsf V,D,(f_c)_{c\in\Lang}),
\]
where \(\mathsf V\) is a finite nonempty set of values,
\(D\subseteq\mathsf V\) is the set of designated values, and every
\(n\)-ary connective \(c\) is interpreted by a total operation
\(f_c:\mathsf V^n\to\mathsf V\).  A valuation
\(v:\mathsf{Prop}\to\mathsf V\) extends uniquely and homomorphically to all
formulas.  The matrix includes its operations and designated set: the number
of values alone does not determine a logic.
\index{logical matrix}
\index{designated value}
\end{hbdefinition}

\begin{hbdefinition}{def:p08-matrix-consequence}{Matrix consequence}
For a formula set \(\Gamma\) and formula \(A\), write
\(\Gamma\entails_{\mathcal M}A\) iff every \(\mathcal M\)-valuation that
assigns a member of \(D\) to every formula in \(\Gamma\) also assigns a member
of \(D\) to \(A\).  Thus matrix consequence preserves designation, not a
particular numerical degree.  Empty or universal designated sets are allowed;
their resulting vacuous cases are part of the definition.
\index{matrix consequence}
\end{hbdefinition}

\begin{hbdefinition}{def:p08-value-labelled-formula}{Value labels and branch compatibility}
For a fixed matrix \(\mathcal M\), a \emph{value-labelled formula}
\(A{:}a\), with \(a\in\mathsf V\), asserts that \(A\) has exactly value
\(a\).  A valuation \(v\) is \emph{compatible with a branch} iff
\(v(A)=a\) for every labelled formula \(A{:}a\) on that branch.  Compatibility
is a metalinguistic relation between a valuation and a tableau branch; it is
stronger than merely designating the formulas and is needed in both directions
of the adequacy proof.
\index{value-labelled formula}
\index{branch compatibility}
\end{hbdefinition}

\begin{hbdefinition}{def:p08-saturated-matrix-branch}{Saturated matrix-tableau branch}
A matrix-tableau branch is \emph{saturated} iff, for every occurrence
\(c(A_1,\ldots,A_n){:}a\) on it, there is a tuple
\((a_1,\ldots,a_n)\in\mathsf V^n\) such that
\(f_c(a_1,\ldots,a_n)=a\) and all labels
\(A_1{:}a_1,\ldots,A_n{:}a_n\) occur on the branch.  This exact witness
property, rather than the informal phrase ``nothing more to do,'' is what the
countervaluation induction uses.  It differs from compatibility: saturation
is syntactic branch data, whereas compatibility refers to an external
valuation.
\index{saturated branch!finite matrix tableau}
\end{hbdefinition}

All root, expansion, marking, and closure clauses live on the
\hyperref[sheet:tableaux-matrix]{finite-matrix tableau sheet}.  In
particular, formula occurrences remain on a branch after expansion and are
marked there; no body paragraph silently supplies an additional rule.

\begin{hbtheorem}{thm:p08-finite-matrix-tableau}{proved=annex:p8.matrix.1}{Finite-matrix tableau adequacy}
Let \(\mathcal M\) be a finite logical matrix, let \(\Gamma\) be a finite set
of formulas, and let \(B\) be a formula.  Every tableau in the root family
specified by the \hyperref[sheet:tableaux-matrix]{finite-matrix tableau sheet}
closes iff
\[
  \Gamma\entails_{\mathcal M}B.
\]
The statement includes termination of full expansion.  It does not cover an
infinite value set, first-order quantifiers, or an infinite premise set.
\end{hbtheorem}

\section{Strong Kleene logic and the Logic of Paradox}

Throughout the K3--LP comparison, the language is generated from propositional
variables by \(\neg,\land,\lor\) and has no nullary truth constants.  This
standing is essential to the theorem-set claims below.

\begin{systementry}{entry:p08-k3}{Strong Kleene logic K3}
  \index{Strong Kleene logic}
  \entryaliases{\(\mathbf K_3\); strong three-valued Kleene logic.}
  \entrylanguage{The fixed constant-free propositional language with
    \(\neg,\land,\lor\).}
  \entrysemantics{The matrix with
    \(\mathsf V=\{0,\tfrac12,1\}\), \(D=\{1\}\),
    \(\neg a=1-a\), \(a\land b=\min(a,b)\), and
    \(a\lor b=\max(a,b)\).}
  \entrycalculus{sheet:tableaux-matrix}
  \entryadequacy{proved=annex:p8.k3.1}{For every finite premise set, the
    registered finite-matrix tableau closes exactly for K3 consequence.}
  \entrydecidability{proved=annex:p8.k3.2}{A three-valued truth-table test
    decides finite-premise K3 consequence.}
  \entrycharacteristic{Excluded middle can fail, and there are no tautologies
    in the fixed constant-free language; see
    Theorem~\ref{thm:p08-k3-lp-facts}.}
  \entrytranslations{It is the FDE submatrix on
    \(\{f,n,t\}\), with only \(t\)
    designated.}
  \entryalgebra{The three-element strong Kleene De Morgan algebra with one
    designated value.}
  \entrysources{Kleene1952,OmoriWansing2017}
\end{systementry}

\begin{systementry}{entry:p08-lp}{Logic of Paradox LP}
  \index{Logic of Paradox}
  \entryaliases{\(\mathsf{LP}\).}
  \entrylanguage{The same fixed constant-free language as K3.}
  \entrysemantics{The same value set and strong Kleene operations as K3, but
    with \(D=\{\tfrac12,1\}\).  The middle value is read as carrying both
    positive and negative support.}
  \entrycalculus{sheet:tableaux-matrix}
  \entryadequacy{proved=annex:p8.lp.1}{For every finite premise set, the
    registered finite-matrix tableau closes exactly for LP consequence.}
  \entrydecidability{proved=annex:p8.lp.2}{A three-valued truth-table test
    decides finite-premise LP consequence.}
  \entrycharacteristic{Explosion fails, while the tautologies are exactly the
    classical tautologies in the fixed language; see
    Theorem~\ref{thm:p08-k3-lp-facts}.}
  \entrytranslations{It is the FDE submatrix on
    \(\{f,b,t\}\), with \(b\) and
    \(t\) designated.}
  \entryalgebra{The same three-element strong Kleene algebra as K3 with a
    different designated filter.}
  \entrysources{Priest1979,Priest2006,OmoriWansing2017}
\end{systementry}

\begin{hbexample}{def:p08-k3-lp-excluded-middle}{One table, two designated sets}
If \(v(p)=\tfrac12\), then
\[
  v(p\lor\neg p)=\max(\tfrac12,\tfrac12)=\tfrac12.
\]
This is nondesignated in K3 and designated in LP.  Hence excluded middle is
not K3-valid but this valuation does not refute it in LP.  The calculation is a
countervaluation/designation check, not a derivation.
\end{hbexample}

\begin{hbexample}{def:p08-lp-explosion}{LP countervaluation to explosion}
Set \(v(p)=\tfrac12\) and \(v(q)=0\).  Both \(p\) and \(\neg p\) are
designated, while \(q\) is not, so
\[
  p,\neg p\not\entails_{\mathsf{LP}}q.
\]
This witnesses failure of unrestricted explosion; it does not assert that a
particular contradiction is derivable.
\end{hbexample}

\begin{hbdefinition}{def:p08-paraconsistent-consequence}{Paraconsistent and paracomplete consequence}
A consequence relation \(\entails_{S}\) is \emph{paraconsistent} iff some
formulas \(A,B\) satisfy
\(A,\neg A\not\entails_{S}B\).  This is a property of a consequence relation,
not a claim that every contradiction is designated or derivable.  The term is
needed to distinguish LP and FDE from K3 and is used in their card fields and
matrix comparison.  It is \emph{paracomplete} iff some formula \(A\) satisfies
\(\varnothing\not\entails_{S}(A\lor\neg A)\).  Thus paraconsistency records failure of
unrestricted explosion from contradictory premises, whereas paracompleteness
records failure of excluded middle as a theorem; neither property implies the
other.  The second distinction is needed for the FDE card and comparison table.
\index{paraconsistent consequence}
\index{paracomplete consequence}
\end{hbdefinition}

\begin{hbwarning}{def:p08-matrix-values-formal}{Truth-value readings are not matrix data}
Words such as ``indeterminate,'' ``both,'' and ``neither'' explain intended
readings.  They do not replace the formal value set, operations, and designated
set.  K3 and LP demonstrate the point sharply: the same numerical truth tables
produce different consequence relations because designation differs.
\end{hbwarning}

\section{First-degree entailment}

\begin{hbdefinition}{def:p08-fde-matrix}{The FDE support-pair matrix}
Let
\[
 \mathrm{BD}=\{(0,1),(0,0),(1,1),(1,0)\}
 =\{f,n,b,t\},
\]
in the displayed order.  For a value \(x=(x^+,x^-)\), the coordinates are its
positive and negative support indicators.  They are needed because a single
linear scale cannot distinguish neither support from both support.  For
\(x=(x^+,x^-)\) and \(y=(y^+,y^-)\), put
\[
\begin{aligned}
 \neg x&=(x^-,x^+),\\
 x\land y&=(x^+\land y^+,\ x^-\lor y^-),\\
 x\lor y&=(x^+\lor y^+,\ x^-\land y^-),
\end{aligned}
\]
where the coordinate operations are Boolean.  Designate
\(\{b,t\}\), exactly the values with positive support.
This finite matrix defines the displayed propositional first-degree-entailment
semantics \cite{Dunn1976,Belnap1977,OmoriWansing2017}.
\index{first-degree entailment}
\index{FDE}
\end{hbdefinition}

\begin{hbdefinition}{def:p08-fde-orders}{Orders on the BD algebra}
For \(x,y\in\mathrm{BD}\), define
\[
\begin{aligned}
 x\leq_t y
 &\quad\text{ iff }\quad
 x^+\leq y^+\ \text{and}\ y^-\leq x^-,\\
 x\leq_k y
 &\quad\text{ iff }\quad
 x^+\leq y^+\ \text{and}\ x^-\leq y^-.
\end{aligned}
\]
The first compares truth support against falsity support; the second compares
total information componentwise.  They are two relations on the same values,
not two consequence relations.  Their nonredundancy is visible at once:
\(n\) and \(b\) are incomparable in the truth order, whereas
they are respectively bottom and top in the information order.  A
\emph{bilattice} is an algebraic structure whose carrier bears two lattice
orders.  The displayed carrier qualifies: truth-order meet and join are
\((\min(x^+,y^+),\max(x^-,y^-))\) and
\((\max(x^+,y^+),\min(x^-,y^-))\), while information-order meet and join take
coordinatewise minima and maxima.  Thus ``bilattice'' records both lattice
orders at once and cannot be replaced by either order alone.  This distinction
is used by the FDE card.
\index{bilattice}
\index{truth order}
\index{information order}
\end{hbdefinition}

\begin{systementry}{entry:p08-fde}{First-degree entailment FDE}
  \index{first-degree entailment!system entry}
  \entryaliases{Belnap--Dunn four-valued logic; \(\mathbf{FDE}\).}
  \entrylanguage{The fixed constant-free propositional language with
    \(\neg,\land,\lor\).}
  \entrysemantics{The support-pair matrix of
    Definition~\ref{def:p08-fde-matrix}, designating
    \(b,t\).}
  \entrycalculus{sheet:tableaux-matrix}
  \entryadequacy{proved=annex:p8.fde.1}{For every finite premise set, the
    registered finite-matrix tableau closes exactly for FDE consequence.}
  \entrydecidability{proved=annex:p8.fde.2}{Finite-premise FDE consequence is
    decidable by exhaustive four-valued evaluation.}
  \entrycharacteristic{It is both paraconsistent and paracomplete: a both-valued
    premise contradiction does not explode, and a neither-valued atom refutes
    excluded middle.}
  \entrytranslations{K3 deletes \(b\) and designates only
    \(t\); LP deletes \(n\) and designates
    \(b,t\).}
  \entryalgebra{\(\mathrm{BD}\) carries the distinct truth and knowledge orders of
    Definition~\ref{def:p08-fde-orders}.}
  \entrysources{Dunn1976,Belnap1977,OmoriWansing2017}
\end{systementry}

\begin{hbtheorem}{thm:p08-k3-lp-facts}{proved=annex:p8.matrix.2}{K3--LP theorem sets and dual consequence}
In the fixed constant-free \(\neg,\land,\lor\)-language:
\begin{enumerate}
  \item \(\mathbf K_3\) has no tautologies;
  \item the LP tautologies are exactly the classical tautologies;
  \item for finite formula sets \(\Gamma,\Delta\), if
        \(\Gamma\entails_S\Delta\) means that every valuation designating all
        members of \(\Gamma\) designates at least one member of \(\Delta\),
        then
        \[
          \Gamma\entails_{\mathbf K_3}\Delta
          \quad\text{ iff }\quad
          \neg\Delta\entails_{\mathsf{LP}}\neg\Gamma,
        \]
        where negation acts memberwise on each finite set.
\end{enumerate}
The third clause is the registered K3--LP duality.  It is not an unqualified
single-conclusion equivalence \cite{Priest1979,OmoriWansing2017}.
\end{hbtheorem}

\begin{hbdefinition}{def:p08-weak-kleene-matrix}{Weak Kleene matrix}
The \emph{weak Kleene matrix} has values
\(\{0,\tfrac12,1\}\), designates only \(1\), fixes
\(\neg\tfrac12=\tfrac12\), and makes \(\tfrac12\) infectious for every
nonunary classical connective: if any immediate argument has that value, so
does the compound; otherwise the operation is classical.  Infectiousness is
an exact operation clause, not the informal claim that a sentence is
meaningless.  This definition is used in Table~\ref{tab:p08-many-valued-matrices}
and nowhere as a substitute for Bochvar's richer external language
\cite{Kleene1952}.
\index{weak Kleene logic}
\end{hbdefinition}

\section{Finite and infinite Lukasiewicz logics}

\begin{proofboundary}{def:p08-many-valued-card-boundary}{Why only three cards appear}
The inherited-card list for this chapter names K3 and LP; the FDE card is
additional later material because the exact finite matrix also uses the real
generic sheet.  Register rows D-P10-020, N-b-028, and N-b-029 remain OPEN
inclusion obligations whose legacy treatment or priority fields do not limit
their required presentation.  The material below records some relevant
content but does not by itself close those rows.  The infinite matrix is not governed by
\texttt{sheet:tableaux-matrix}, and this visit licenses no infinite-valued
Hilbert sheet; printing a \texttt{systementry} card would therefore fabricate
its compulsory calculus field.  Thus every printed card has a real exact
sheet, while none of the register rows is silently advertised as closed by a
card, definition, theorem, table, or pointer.
\index{proof boundary!many-valued system cards}
\end{proofboundary}

\begin{hbdefinition}{def:p08-lukasiewicz-matrices}{Cignoli's finite chain and H\'ajek's standard \(\text{\L}\) semantics}
For \(n\geq2\), Cignoli, D'Ottaviano, and Mundici write the finite
MV-chain itself as
\[
 \text{\L}_n=\left\{0,\frac1{n-1},\frac2{n-1},\ldots,1\right\}.
\]
The \(n\)-valued Lukasiewicz system uses this value algebra, designated set
\(D=\{1\}\), and the operations
\[
  \neg x=1-x,
  \qquad
  x\Rightarrow y=\min(1,1-x+y).
\]
H\'ajek writes the standard infinite-valued logic as \(\text{\L}\); its
standard value algebra is \([0,1]\).  Its derived strong conjunction is
\(x\mathbin{\ast_{\text{\L}}}y=\max(0,x+y-1)\); weak conjunction and
disjunction are minimum and maximum.  These are H\'ajek's source symbols:
\(\Rightarrow\) is typed here as the truth-function implication and
\(\ast_{\text{\L}}\) as the Lukasiewicz t-norm.  A finite system based on
\(\text{\L}_n\) uses
the generic finite-matrix sheet; the standard \(\text{\L}\) semantics does not
\cite{Lukasiewicz1920,LukasiewiczTarski1930,Hajek1998,
  CignoliDOttavianoMundici2000}.
\index{Lukasiewicz logic}
\index{many-valued logic!Lukasiewicz}
\end{hbdefinition}

\begin{landscape}
\begin{longtable}{@{}p{.10\linewidth}p{.16\linewidth}p{.13\linewidth}p{.27\linewidth}p{.25\linewidth}@{}}
\caption{Many-valued matrix comparison (apparatus A-084).  The theorem-set
column uses the fixed constant-free language for K3, LP, weak Kleene, and FDE.
The finite Lukasiewicz row uses its separately displayed implication language.}
\label{tab:p08-many-valued-matrices}\\
\toprule
System & Values & Designated & Operations & Theorem set and boundary \\
\midrule
\endfirsthead
\toprule
System & Values & Designated & Operations & Theorem set and boundary \\
\midrule
\endhead
\(\mathbf K_3\) & \(0,\tfrac12,1\) & \(1\) & Strong Kleene:
negation \(1-a\), conjunction minimum, disjunction maximum. & Empty theorem
set in the fixed language.  The middle value is nondesignated; finite
multiple-conclusion dual to LP. \\
\addlinespace
\(\mathsf{LP}\) & \(0,\tfrac12,1\) & \(\tfrac12,1\) & The same strong
Kleene operations. & Exactly the classical tautologies, but not classical
consequence: explosion fails.  The middle value is designated. \\
\addlinespace
Weak Kleene matrix & \(0,\tfrac12,1\) & \(1\) &
The middle value is infectious; otherwise operations are classical. & Empty
theorem set in the fixed language.  This is the internal matrix row, not a
card for full Bochvar external logic. \\
\addlinespace
\(\mathbf{FDE}\) & \(f,n,b,t\) &
\(b,t\) & Support-pair De Morgan operations from
Definition~\ref{def:p08-fde-matrix}. & Empty theorem set in the fixed language;
both paraconsistent and paracomplete.  K3 and LP are distinct three-value
submatrices. \\
\addlinespace
The \(n\)-valued \(\text{\L}\) system, with value chain \(\text{\L}_n\),
\(n\geq2\) &
\(0,\frac1{n-1},\ldots,1\) & \(1\) & Lukasiewicz negation and implication
from Definition~\ref{def:p08-lukasiewicz-matrices}. & The finite-valued
Lukasiewicz theorem set.  \(n=2\) is classical; every finite row is decided by
\texttt{sheet:tableaux-matrix}. \\
\bottomrule
\end{longtable}
\end{landscape}

The K3 and LP tableaux share the same expansion shapes because their
operations agree, but their root families differ because their designated sets
differ.  This is the reusable content of the source's first exercise; it does
not identify the two consequence relations.

\hbstub{D-P10-025 remains open at its could ceiling: Table A-084 supplies the
weak Kleene internal matrix, but Bochvar's full external language, assertion
operators, exact calculus, and adequacy theorem have not been supplied from one
cited source presentation.}

\hbstub{D-P10-023 is a could-priority row: no particular family of Post
many-valued logics has yet been supplied with its source-defined language,
matrices, and
status-bearing theorem, so a generic name is not presented as a result.}

\begin{hbtheorem}{thm:p08-lukasiewicz-completeness}{cited=Hajek1998}{Lukasiewicz completeness}
For its standard Hilbert calculus, in the language of
Definition~\ref{def:p08-lukasiewicz-matrices}:
\begin{enumerate}
  \item a formula is a theorem iff every valuation in the standard
        \([0,1]\)-matrix gives it value \(1\);
  \item for a finite premise set, syntactic consequence agrees with preservation
        of value \(1\) in that standard matrix.
\end{enumerate}
These clauses do not license arbitrary-premise strong completeness for the
single standard matrix.  The exact definition of the corresponding algebraic
semantics and its arbitrary-premise completeness theorem are deferred to
registered rows D-P10-022 and N-d-013; no algebra class is silently introduced
here \cite{CignoliDOttavianoMundici2000}.
\end{hbtheorem}

\begin{hbtheorem}{thm:p08-mcnaughton}{cited=McNaughton1951}{McNaughton's theorem}
A function \(g:[0,1]^m\to[0,1]\) is the truth function of an
\(m\)-variable infinite-valued Lukasiewicz formula iff \(g\) is continuous and
piecewise linear with integer coefficients on every linear piece.  The theorem
does not represent arbitrary continuous functions, and it is not an
arbitrary-premise completeness statement.
\end{hbtheorem}

For each finite \(n\), Theorem~\ref{thm:p08-finite-matrix-tableau} supplies a
terminating decision procedure for consequence in a matrix based on
\(\text{\L}_n\).  The cited
Hilbert axiomatization concerns the infinite logic; no finite-tableau claim is
extended to the real interval, and the deferred algebraic theorem is not used
to strengthen that claim.

\section{T-norm fuzzy logics}

\begin{hbdefinition}{def:p08-tnorm-residuum}{T-norm and residuum}
A \emph{t-norm} is a commutative, associative, coordinatewise monotone
operation \(\ast:[0,1]^2\to[0,1]\) with unit \(1\).  If \(\ast\) is
left-continuous, its \emph{residuum} is
\[
  x\Rightarrow y
  =\sup\{z\in[0,1]:x\ast z\leq y\}.
\]
The t-norm interprets strong conjunction, whereas lattice conjunction remains
minimum; the residuum is a semantic operation interpreting object-language
implication.  The source uses the plain glyph \(\Rightarrow\); its numeric
operands type it as a residuum rather than a sequent arrow.  These distinctions
are needed in every row of Table~\ref{tab:p08-tnorm-logics}
\cite{Hajek1998,CintulaHajekNoguera2011}.
\index{t-norm}
\index{residuum!t-norm}
\end{hbdefinition}

\begin{hbtheorem}{thm:p08-tnorm-completeness}{cited=CintulaHajekNoguera2011}{Propositional t-norm completeness survey}
The calculi named in Table~\ref{tab:p08-tnorm-logics} have the cited standard
completeness surfaces stated there.  In particular, MTL ranges over
left-continuous t-norms and BL over continuous t-norms.  No theoremhood or
finite-premise standard result in the table is promoted to arbitrary-premise
strong completeness for a single standard chain.  Exact algebra-class
definitions and their completeness theorems are deferred to registered row
N-d-015 in Part 10.
\end{hbtheorem}

\begin{landscape}
\begin{longtable}{@{}p{.10\linewidth}p{.18\linewidth}p{.20\linewidth}p{.20\linewidth}p{.23\linewidth}@{}}
\caption{T-norm fuzzy logics (apparatus A-085).  ``Axiomatization'' names the
exact cited Hilbert extension without redisplaying rules; none of these rows
asserts that a P8.5 Hilbert sheet exists.}
\label{tab:p08-tnorm-logics}\\
\toprule
Logic & Standard t-norm & Residuum & Cited calculus/semantic scope & Completeness surface \\
\midrule
\endfirsthead
\toprule
Logic & Standard t-norm & Residuum & Cited calculus/semantic scope & Completeness surface \\
\midrule
\endhead
\(\mathrm{MTL}\) & Every left-continuous t-norm. &
\(\sup\{z:x\ast z\leq y\}\). & Esteva--Godo MTL; standard semantics ranges
over all left-continuous t-norms. & Standard theoremhood completeness over all
left-continuous t-norms \cite{EstevaGodo2001,JeneiMontagna2002}; the exact
algebraic counterpart is deferred to N-d-015. \\
\addlinespace
\(\mathrm{BL}\) & Every continuous t-norm. & The displayed residuum. &
H\'ajek's cited BL calculus; standard semantics ranges over all continuous
t-norms. & Standard theoremhood and finite-premise completeness
over all continuous t-norms \cite{CignoliEstevaGodoTorrens2000}; the exact
algebraic counterpart is deferred to N-d-015. \\
\addlinespace
\(\text{\L}\) & \(\max(0,x+y-1)\). &
\(\min(1,1-x+y)\). & The cited Lukasiewicz extension of BL, with the standard
operations fixed in Definition~\ref{def:p08-lukasiewicz-matrices}. & Standard
theoremhood and finite-premise completeness
\cite{Hajek1998,CignoliDOttavianoMundici2000}; the arbitrary-premise algebraic
theorem is deferred to D-P10-022 and N-d-013. \\
\addlinespace
\(G\) & \(\min(x,y)\). & \(1\) if \(x\leq y\), and
\(y\) otherwise. & The cited G\"odel extension of BL; its standard conjunction
satisfies \(x\ast x=x\). & Standard theoremhood completeness for the
G\"odel chain \cite{Hajek1998}; the exact algebraic counterpart is deferred
to N-d-015. \\
\addlinespace
\(\Pi\) & \(x\mathbin{\ast_{\Pi}}y=xy\). & \(1\) if \(x\leq y\), and \(y/x\) otherwise.
& H\'ajek's cited two-schema Product extension of BL, with multiplication as
its standard t-norm. & Standard theoremhood completeness for the Product chain
\cite{Hajek1998}; the exact algebraic counterpart is deferred to N-d-015. \\
\bottomrule
\end{longtable}
\end{landscape}

The table's continuity column is part of the mathematical statement: BL is not
the logic of all left-continuous t-norms, and MTL is not restricted to the
continuous ones.  Likewise, the three named standard chains are examples with
different strong conjunctions, not alternative informal readings of one fixed
matrix.

\section{Supervaluations}

\begin{hbdefinition}{def:p08-supervaluation}{Fine specification spaces and supertruth}
Following Fine's presentation, a \emph{specification space} is a partially
ordered collection of specification-points, where the order records extension
of information.  A \emph{base-point} is the specification-point being
evaluated.  A point above it is a \emph{completion} when it gives a complete
classical specification, and it is \emph{admissible} when it satisfies the
primitive admissibility constraints of that specification space.  The
base-point is \emph{completeable} when at least one complete admissible point
extends it.

At a completeable base-point, a formula is \emph{supertrue} iff it is true at
every complete admissible specification extending that point,
\emph{superfalse} iff it is false at every such specification, and
\emph{indeterminate} otherwise.  Specification-points and admissibility are
metalinguistic semantic data, not matrix truth values.  The extension order is
needed to keep only completions of the base information, while admissibility
can further restrict those completions; neither condition subsumes the other.
This distinction is used in the preservation theorem and example below
\cite{Fine1975}.
\index{supervaluation}
\index{admissible precisification}
\index{supertruth}
\end{hbdefinition}

\begin{hbtheorem}{thm:p08-supervaluation-classical}{cited=Fine1975}{Classical tautologies under supervaluation}
Every classical tautology is supertrue at every completeable base-point.  The
converse holds when every complete classical specification is an admissible
completion of the base-point, but need not hold when primitive admissibility
restricts its completions.
\end{hbtheorem}

\begin{hbexample}{def:p08-supervaluation-gap}{Supertruth with an indeterminate atom}
Suppose two complete admissible specifications extending one base-point
disagree about \(p\).  Then \(p\) is neither supertrue nor superfalse, but every
such specification makes
\(p\lor\neg p\) true, so that compound is supertrue.  Supertruth therefore is
not a truth-functional third value assigned compositionally from the three
supervaluation statuses of immediate constituents.
\end{hbexample}

\section{Sources and navigation for this chapter}

The inherited matrix definitions, exact generic sheet, repaired theorem,
K3/LP matrices, two countervaluations, and formal-data warning migrate v1.1
lines 13584--13787.  The inherited proof is
Annex~\ref{annex:p8.matrix.1}.  The K3, LP, and FDE cards all point to the one
real generic sheet; their finite specialization and decision arguments have
distinct annex targets.  Finite Lukasiewicz matrices point to that sheet in
prose, while the infinite real matrix expressly does not.

The new FDE entry follows \cite{Dunn1976,Belnap1977,OmoriWansing2017}; the
Lukasiewicz entry follows
\cite{LukasiewiczTarski1930,McNaughton1951,CignoliDOttavianoMundici2000}; and
the t-norm survey follows
\cite{Hajek1998,EstevaGodo2001,CintulaHajekNoguera2011}.  These citations and
the matrix/t-norm notation are this chapter's contributions to apparatus
A-165, A-182, and A-189.  The single Part-8 further-reading section remains a
part-finalizer responsibility; this chapter does not claim A-186.
